\chapter{题型详细判断树}

\section{本章定位}

本章是整本手册的导航核心。考场上不要从算法目录开始翻，而应先用本章判断：

\begin{enumerate}
  \item 这题属于哪一类？
  \item 数据范围允许什么复杂度？
  \item 题面信号对应哪个模板？
  \item 如果满分不会，有什么部分分做法？
\end{enumerate}

\Warn{判断树不是保证一眼看出满分算法的魔法。}它的作用是减少迷路：让你从题面信号稳定定位到最可能的算法和模板。

\section{考场三层分流法}

读任何题都按三层走。不要直接猜算法名。

\begin{longtable}{p{3cm}p{5.2cm}p{6.2cm}}
\toprule
层级 & 问什么 & 作用 \\
\midrule
第一层：题号和任务 & A/B/C/D/E？要求输出一个值、一组答案，还是处理很多命令？ & 决定目标：稳拿、冲分、拆模拟、拿部分分。 \\
第二层：题面结构 & 输入是数组、矩阵、图、字符串、对象、事件、树，还是混合规则？ & 决定先翻哪个章节。 \\
第三层：操作和范围 & 是否多次查询/修改？$n,m,q$ 多大？值域多大？ & 决定具体模板和复杂度。 \\
\bottomrule
\end{longtable}

\section{总判断流程}

\begin{longtable}{p{2.6cm}p{6.2cm}p{5.8cm}}
\toprule
步骤 & 问题 & 动作 \\
\midrule
1 & 这是第几题？ & A/B 先求稳，C 先拆题，D/E 先找部分分。 \\
2 & 输出形式是什么？ & 单个答案偏统计/公式；多行答案偏查询/命令；复杂文本偏解析。 \\
3 & 数据范围多大？ & 用复杂度表排除暴力或高级算法。 \\
4 & 有多次查询或修改吗？ & 查前缀和、差分、树状数组、线段树、离线处理。 \\
5 & 有图、网格、路径吗？ & 查 BFS、DFS、Dijkstra、并查集、拓扑排序。 \\
6 & 有字符串格式吗？ & 查 split、路径、URL、括号、表达式、KMP。 \\
7 & 有状态、命令、事件吗？ & 查大模拟、状态机、事件排序。 \\
8 & 有最优值、选择、容量吗？ & 查 DP、贪心、二分答案。 \\
9 & 满分算法不明确吗？ & 立即找小数据和特殊性质分。 \\
\bottomrule
\end{longtable}

\section{决策树叶子节点标准格式}

本章后面的详细决策树尽量使用同一种叶子格式。考场上每次走到叶子，都要确认 5 件事：

\begin{longtable}{p{2.8cm}p{11.2cm}}
\toprule
项目 & 必须回答的问题 \\
\midrule
题面信号 & 是什么词、输入结构或输出要求让我走到这里？ \\
数据范围 & 暴力能不能过？目标复杂度应是多少？ \\
模板入口 & 应翻第几章、哪个模板编号？ \\
易错点 & 下标、边界、溢出、排序规则、输出格式哪里最容易错？ \\
保底做法 & 满分不会时，能不能先写暴力、小数据或特殊性质版本？ \\
\bottomrule
\end{longtable}

\Warn{如果某个叶子节点无法回答“数据范围”和“模板入口”，不要立刻写代码，先回到数据范围判断树重新分流。}

\section{从根节点开始的总决策树}

下面这棵树按考场读题顺序走。先用它确定大方向，再进入 A/B/C/D/E 专项树。

\begingroup
\small
\setlength{\tabcolsep}{3pt}
\begin{longtable}{p{1.5cm}p{4.0cm}p{4.9cm}p{4.1cm}}
\toprule
节点 & 问题 & 是：走向 & 否：走向 \\
\midrule
R0 & 这是 A 题吗？ & 进入 A 题详细树，先求稳。 & 看 R1。 \\
R1 & 这是 B 题吗？ & 进入 B 题详细树，先找重复计算。 & 看 R2。 \\
R2 & 这是 C 题吗？ & 进入 C 题详细树，先建模。 & 看 R3。 \\
R3 & 这是 D/E 题吗？ & 先看子任务，进入部分分树。 & 回题面确认题号。 \\
R4 & 输出只有一个数或少量值吗？ & 优先统计、公式、排序、DP。 & 看是否多次查询或命令。 \\
R5 & 每次输入操作后都要输出吗？ & 命令分发、数据结构、状态机。 & 可能是最后统一输出的模拟/预处理。 \\
R6 & 有 $q$ 次查询/修改且 $q$ 大吗？ & 前缀、差分、树状数组、线段树、离线。 & 暴力或简单模拟可能可行。 \\
R7 & 有“最小/最大/阈值/最少资源”吗？ & 检查单调性，可能二分答案或 DP。 & 看图、字符串、状态。 \\
R8 & 有“路径/连通/最短/依赖”吗？ & 图论或拓扑。 & 看字符串和大模拟。 \\
R9 & 有复杂格式、路径、括号、表达式吗？ & 字符串解析、栈、递归解析。 & 普通数组/矩阵/对象模拟。 \\
\bottomrule
\end{longtable}
\endgroup

\section{详细叶子索引总表}

这张表是全局的“最终落点”索引。专项树走不清楚时，直接查这一张。

\begingroup
\small
\setlength{\tabcolsep}{3pt}
\begin{longtable}{p{2.8cm}p{3.2cm}p{3.1cm}p{3.0cm}p{2.3cm}}
\toprule
叶子 & 题面信号 & 数据范围判断 & 模板入口 & 保底做法 \\
\midrule
L01 统计扫描 & 求和、计数、最大最小 & $O(n)$ 可过 & A 题、T106 & 一遍循环。 \\
L02 排序排名 & 排名、中位数、多关键字 & $O(n\log n)$ 可过 & T105 & 先按单关键字过样例。 \\
L03 一维区间 & 多次区间和/数量 & 静态查询多 & T101 & 每次暴力扫小数据。 \\
L04 二维矩阵 & 子矩阵、邻域、图像 & $nm$ 可扫，多次查询用预处理 & T102/T104 & 小矩阵暴力。 \\
L05 哈希查重 & 是否出现、频率、点集 & 值域大或字符串键 & T106 & map 替代 unordered\_map。 \\
L06 坐标离散 & 坐标大、点少 & 不能直接开数组 & T109 & map 坐标。 \\
L07 二分答案 & 最小时间、最大阈值 & check 单调且快 & T107 & 枚举小范围答案。 \\
L08 背包选择 & 选若干、容量、金额 & 总容量可 DP & T301/T302 & DFS 小数据。 \\
L09 BFS & 最短步数、边权 1 & 状态数可控 & T201 & 小图暴力搜索。 \\
L10 Dijkstra & 正权最短路 & $n,m$ 大，边权非 1 & T204 & 小图 Floyd。 \\
L11 拓扑 & 前置、依赖、逻辑门 & DAG 或需判环 & T205 & 小图反复找入度 0。 \\
L12 并查集 & 合并集合、连通性 & 只合并查连通 & T203 & BFS/DFS 小数据。 \\
L13 事件模拟 & 按时间发生、到期 & 事件数可排序 & T110 & 按输入顺序模拟小数据。 \\
L14 状态机 & 对象状态变化 & 状态少且规则明确 & T603 & if/else 先过样例。 \\
L15 路径/URL & 路径、参数、配置 & 总长度可线性 & T606/T607 & 手写 split。 \\
L16 表达式 & 括号、优先级、化学式 & 嵌套深度可递归 & T608/T609 & 暴力解析小表达式。 \\
L17 树上路径 & 祖先、子树、路径 & 查询多需预处理 & T208/T402 & 每次 DFS。 \\
L18 动态区间 & 修改和查询交替 & $q$ 大 & T401/T402 & 修改少就重算。 \\
L19 大模拟 & 长题面、多对象规则 & 规则多于算法 & T601/T602 & 分命令逐步拿分。 \\
L20 D/E 小数据 & 子任务 $n$ 小或特殊性质 & 暴力能过一档 & 第 8 章 & 先交可运行版本。 \\
\bottomrule
\end{longtable}
\endgroup

\section{题面结构第一判断}

\begingroup
\small
\setlength{\tabcolsep}{3pt}
\begin{longtable}{p{3.1cm}p{4.1cm}p{4.2cm}p{3.0cm}}
\toprule
输入主体 & 题面信号 & 第一反应 & 翻到哪里 \\
\midrule
一串数字 / 数组 & 求和、最大、最小、相邻、分段、前缀 & 扫描、排序、前缀和 & A/B、T101、T105 \\
二维矩阵 & 图像、网格、矩形、邻域、旋转、重塑 & 下标转换、二维前缀、矩阵模拟 & A/B、T102、T104 \\
很多字符串行 & 路径、URL、JSON、Markdown、命令、表达式 & 字符串解析、大模拟 & C、T501、T606、T607 \\
点和边 & 连通、最短路、依赖、拓扑、最小代价 & 图论模板 & T201--T207 \\
树 & 子树、祖先、路径、删点、合并 & DFS、LCA、树形 DP、DFS 序 & C/D/E、T208 \\
多条操作 & 修改、查询、插入、删除、回滚 & 数据结构或离线处理 & T103、T401、T402 \\
时间顺序事件 & 某时刻发生、到期、优先级、调度 & 事件排序、状态机、堆 & C、T110、T603、T604 \\
选择若干物品 & 容量、费用、价值、至少/至多 & 背包 DP 或贪心 & T301、T302 \\
答案最优值 & 最小时间、最大阈值、最小最大 & 二分答案或 DP & T107、D/E \\
\bottomrule
\end{longtable}
\endgroup

\section{数据范围判断树}

\subsection{单变量规模}

\begin{longtable}{p{3.4cm}p{5.4cm}p{5.4cm}}
\toprule
范围 & 可考虑 & 不宜考虑 \\
\midrule
$n \le 10$ & 全排列、爆搜、手写模拟所有情况 & 不必急着写复杂优化 \\
$n \le 20$ & 子集枚举、状压 DP、折半搜索 & $O(n!)$ 仍危险 \\
$n \le 100$ & $O(n^3)$、Floyd、区间 DP & 指数算法 \\
$n \le 1000$ & $O(n^2)$、简单 DP、稠密图 & $O(n^3)$ 需谨慎 \\
$n \le 10^5$ & $O(n\log n)$、$O(n)$ & $O(n^2)$ \\
$n \le 10^6$ & $O(n)$、计数数组、简单筛法 & 多层 map 或重型结构 \\
$n \ge 10^7$ & 只考虑线性或数学公式 & 大数组套多维结构 \\
\bottomrule
\end{longtable}

\subsection{多变量规模}

\begin{longtable}{p{4.2cm}p{5.2cm}p{5.0cm}}
\toprule
范围组合 & 判断 & 常见模板 \\
\midrule
$n,m \le 500$ 且矩阵 & $O(nm)$ 或 $O(nm\log nm)$ 通常可行 & 二维扫描、T102、T104 \\
$n,m \le 2000$ 且矩阵 & 尽量 $O(nm)$，少做多次全图扫描 & 二维前缀、BFS \\
$q \le 10^5$ & 每次操作必须快 & 前缀、差分、树状数组、线段树 \\
值域 $\le 10^6$ & 可以开计数数组 & 计数、桶、前缀计数 \\
值域很大但数量少 & 不可直接开值域数组 & map、unordered\_map、离散化 \\
边数 $m \approx n$ & 稀疏图 & 邻接表、BFS、Dijkstra 堆优化 \\
边数 $m \approx n^2$ 且 $n$ 小 & 稠密图 & Floyd、邻接矩阵 \\
字符串总长 $\le 10^6$ & 按总长度线性处理 & split、KMP、Trie \\
\bottomrule
\end{longtable}

\Warn{数据范围优先级高于题面关键词。}例如题面写“最短路”，但 $n \le 100$ 且多源查询，可以用 Floyd；若 $n,m \le 10^5$，通常要 Dijkstra 或 BFS。

\section{关键词到模板总索引}

\begingroup
\small
\setlength{\tabcolsep}{3pt}
\begin{longtable}{p{3.5cm}p{3.8cm}p{4.2cm}p{2.9cm}}
\toprule
题面关键词 & 判断 & 模板 & 章节 \\
\midrule
多次区间和 & 静态区间查询 & T101 一维前缀和 & B / 模板库 \\
子矩阵和、邻域均值 & 二维静态查询 & T102 二维前缀和 & B / 模板库 \\
区间加，最后查询 & 离线修改 & T103 一维差分 & B / 模板库 \\
矩形加，最后输出 & 二维离线修改 & T104 二维差分 & B / 模板库 \\
排名、优先级、多关键字 & 比较器排序 & T105 排序比较器 & A/B/C \\
出现次数、频率、查重 & 统计 & T106 哈希统计 & A/B \\
最小化最大值、最大化最小值 & 二分答案 & T107 二分答案 & B/D \\
排序后找两数、合并稀疏向量 & 双指针 & T108 双指针 & B \\
坐标很大、点不多 & 离散化 & T109 坐标离散化 & B/D \\
按时间发生、到期、归还 & 事件模拟 & T110 事件排序 & B/C \\
最短步数、迷宫、扩散 & BFS & T201 网格 BFS & B/C \\
连通块、区域大小 & DFS/BFS & T202 DFS 连通块 & B/C \\
合并集合、同组、连通性 & 并查集 & T203 并查集 & B/D \\
正权最短路 & Dijkstra & T204 Dijkstra & B/D \\
依赖、先后顺序、任务计划 & 拓扑排序 & T205 拓扑排序 & B/D \\
任意两点最短路，小图 & Floyd & T206 Floyd & B/D \\
最小连接代价 & Kruskal & T207 Kruskal & D/E \\
树上祖先、路径 & LCA & T208 LCA & D/E \\
选或不选、容量 & 01 背包 & T301 01 背包 & B/D \\
可重复选择 & 完全背包 & T302 完全背包 & B/D \\
单点修改区间和 & 树状数组 & T401 树状数组 & D/E \\
区间修改查询 & 线段树 & T402 线段树 & D/E \\
最近更大/更小 & 单调栈 & T403 单调栈 & B/D \\
滑动窗口极值 & 单调队列 & T404 单调队列 & B/D \\
路径、配置分割 & split & T501 split & C \\
长串匹配 & KMP & T502 KMP & D/E \\
前缀查询 & Trie & T503 Trie & D/E \\
分组、分页、负数取模 & 整除取整 & T510 整除取整和标准取模 & A/B \\
快速幂、取模幂 & 数学 & T511 快速幂 & B/D \\
素数、约数、因子 & 数学 & T512 素数约数 & B/D \\
进制转换、编码 & 数字处理 & T513 进制转换 & A/B/C \\
日期、时间、周期 & 日期处理 & T514 日期处理 & A/C \\
长题面、多命令 & 大模拟 & T601 大模拟骨架 & C \\
每行操作名 & 命令分发 & T602 命令分发 & C \\
多状态对象 & 状态机 & T603 状态机 & C \\
动态删除仍取最优 & 懒删除堆 & T604 懒删除优先队列 & C/D \\
最近最少使用 & LRU & T605 LRU 缓存 & C \\
目录路径 & 路径规范化 & T606 路径规范化 & C \\
URL 参数 & URL 解析 & T607 URL 参数解析 & C \\
括号合法性 & 栈 & T608 括号匹配 & C \\
简单表达式 & 扫描求值 & T609 表达式入门 & C \\
\bottomrule
\end{longtable}
\endgroup

\section{操作模式判断树}

\subsection{只有查询，没有修改}

\begin{longtable}{p{4.6cm}p{5.4cm}p{4.0cm}}
\toprule
查询内容 & 优先算法 & 模板 \\
\midrule
区间和、区间数量 & 前缀和 / 前缀计数 & T101 \\
子矩阵和、邻域平均 & 二维前缀和 & T102 \\
最大最小，静态 & 排序 / 预处理 / 单调结构 & T105 / T403 \\
图上单源距离 & BFS / Dijkstra & T201 / T204 \\
任意两点距离，小图 & Floyd & T206 \\
字符串是否出现 & find / KMP & T502 \\
子树信息，静态 & DFS 序 + 前缀 / 树状数组 & T208 / T401 \\
\bottomrule
\end{longtable}

\subsection{有修改，有查询}

\begin{longtable}{p{4.6cm}p{5.4cm}p{4.0cm}}
\toprule
操作形态 & 优先算法 & 模板 \\
\midrule
单点修改，区间和查询 & 树状数组 & T401 \\
区间修改，单点查询 & 差分思想 / 树状数组 & T103 / T401 \\
区间修改，区间查询 & 线段树 & T402 \\
动态插入删除，查最大最小 & set / 堆 / 懒删除 & T604 \\
修改很少 & 每次重算 & 部分分策略 \\
查询很少 & 每次暴力 & 部分分策略 \\
操作可以离线 & 排序事件、差分、并查集倒序 & T110 / T103 / T203 \\
\bottomrule
\end{longtable}

\subsection{多轮模拟}

\begin{longtable}{p{4.6cm}p{5.4cm}p{4.0cm}}
\toprule
模拟形态 & 判断重点 & 模板 \\
\midrule
每轮所有对象同时变化 & 先存变化量，再统一更新 & T601 / T603 \\
事件按时间发生 & 先排序，时间相同按题目优先级 & T110 \\
状态会过期 & 每次操作前处理过期，或懒处理 & T603 / T604 \\
每次选最优对象 & priority\_queue / set & T604 \\
命令种类多 & 每种命令写独立函数 & T602 \\
\bottomrule
\end{longtable}

\section{A 题快速判断树}

A 题目标是稳。判断时按下面顺序，不要一开始就想复杂算法。

\begin{enumerate}
  \item 是否只要求一个数？优先一遍扫描、公式、计数。
  \item 是否有排序、排名、最大最小？用 \Code{sort} 后处理。
  \item 是否有字符或字符串？按字符扫描，先别写复杂解析。
  \item 是否有二维数组？先画 $2 \times 3$ 小例子确认下标。
  \item 是否有日期、时间、进制、取整、小数？翻对应模板，不临场硬想语法。
\end{enumerate}

\begin{longtable}{p{4.6cm}p{5.4cm}p{4.0cm}}
\toprule
题面信号 & 做法 & 模板/章节 \\
\midrule
求和、计数、最大最小 & 一遍扫描 & A 题统计 \\
出现次数、直方图 & 计数数组 / map & T106 \\
小数、百分比、平均、方差 & 公式 + \Code{setprecision} & A 题公式 \\
分组、分页、向上取整、负数取模 & 整数公式，不用浮点硬凑 & T510 \\
排序、排名、中位数 & sort & T105 \\
字符串、字符、校验码 & string 扫描 & T501 \\
进制、数位、编码 & 进制转换或取模拆位 & T513 \\
日期、星期、红绿灯周期 & 闰年、月份、取模 & T514 \\
坐标距离、矩形面积 & 公式，注意边界 & A 题坐标 \\
二维矩阵变换 & 写原坐标到新坐标公式 & A 题二维表格 \\
\bottomrule
\end{longtable}

\subsection{A 题详细叶子表}

\begingroup
\small
\setlength{\tabcolsep}{3pt}
\begin{longtable}{p{2.8cm}p{3.0cm}p{3.0cm}p{3.0cm}p{2.6cm}}
\toprule
叶子 & 识别信号 & 先写什么 & 易错点 & 翻到哪里 \\
\midrule
A1 纯统计 & 求总和、数量、最大最小 & 一遍扫描 & 初值、long long & 第 5 章统计 \\
A2 频率统计 & 出现次数、直方图 & 数组计数或 map & 值域是否能开数组 & T106 \\
A3 排序输出 & 第几大、排名、中位数 & sort 后取值 & 下标从 0 还是 1 & T105 \\
A4 公式计算 & 平均、方差、加权、现值、分组 & 按公式翻译 & 整数除法、取整、小数位 & 第 5 章公式 / T510 \\
A5 字符处理 & 校验码、字符类别、字符串计数 & string 扫描 & getline 和 cin 混用 & T501 \\
A6 数位/进制 & 数位和、编码、进制分解 & 取模除法 & 0、负数、long long & T513 \\
A7 日期时间 & 星期、天数、周期灯 & 月份表 + 闰年 & 闭区间、取模 & T514 \\
A8 二维表格 & 图像、矩阵、邻居平均 & 画小矩阵找下标 & 行列反、边界 & 第 5 章二维 \\
A9 坐标几何 & 距离、矩形面积、点位置 & 套公式 & 交集为负取 0 & 第 5 章坐标 \\
A10 简单状态 & 连续得分、开关、当前状态 & 变量记录状态 & 状态重置时机 & 第 5 章模拟 \\
\bottomrule
\end{longtable}
\endgroup

\section{B 题快速判断树}

B 题目标是尽量拿满。它常常不是“高级算法”，而是把暴力换成预处理。

\begin{enumerate}
  \item 暴力做法是什么？每次是否重复算同一段、同一块、同一类？
  \item 重复算区间：前缀和。重复改区间：差分。
  \item 重复查存在：set/map/hash。
  \item 排序后规律明显：排序扫描、双指针、二分。
  \item 有最小时间/最大阈值：先问答案是否单调。
\end{enumerate}

\begin{longtable}{p{4.6cm}p{5.4cm}p{4.0cm}}
\toprule
题面信号 & 做法 & 模板 \\
\midrule
多次区间/矩阵查询 & 前缀和 & T101 / T102 \\
区间修改最后输出 & 差分 & T103 / T104 \\
频率统计、查重、点集匹配 & 哈希 / map / set & T106 \\
排序后扫描、阈值统计 & sort + 扫描 & T105 \\
两个有序序列匹配 & 双指针 & T108 \\
答案具有单调性 & 二分答案 & T107 \\
最短步数、扩散 & BFS & T201 \\
连通关系、合并集合 & 并查集 & T203 \\
简单依赖、任务计划 & 拓扑排序 & T205 \\
背包信号、至少达到某值 & DP & T301 / T302 \\
质因数、约数、幂 & 数学模板 & T511 / T512 \\
\bottomrule
\end{longtable}

\subsection{B 题详细叶子表}

\begingroup
\small
\setlength{\tabcolsep}{3pt}
\begin{longtable}{p{2.6cm}p{3.2cm}p{3.2cm}p{3.0cm}p{2.6cm}}
\toprule
叶子 & 暴力形态 & 优化判断 & 模板 & 保底 \\
\midrule
B1 一维前缀 & 每次扫区间 & 只查不改，多次区间和 & T101 & 暴力扫小数据 \\
B2 二维前缀 & 每次扫矩形 & 矩阵静态，多次子矩阵 & T102 & 暴力双循环 \\
B3 一维差分 & 每次区间逐个加 & 区间修改，最后输出/查询少 & T103 & 直接修改小数据 \\
B4 二维差分 & 每次矩形逐格加 & 矩形覆盖，最后统计 & T104 & 二维暴力 \\
B5 哈希统计 & 每次线性找存在 & 查重、频率、坐标点存在 & T106 & map 先写稳 \\
B6 排序扫描 & 枚举每个阈值再统计 & 排序后相邻段贡献可更新 & T105/T101 & 枚举阈值小数据 \\
B7 双指针 & 双重循环匹配 & 两序列有序或可排序 & T108 & map 匹配 \\
B8 二分答案 & 枚举答案 & check 单调，答案范围大 & T107 & 枚举小范围 \\
B9 简单 BFS & 枚举路径 & 无权最短步数/扩散 & T201 & DFS 小图 \\
B10 并查集 & 每次 DFS 判连通 & 只有合并和查连通 & T203 & 邻接表 DFS \\
B11 拓扑 & 反复扫描找可做任务 & 前置依赖固定 & T205 & 小 n 反复扫描 \\
B12 背包 & 枚举所有选择 & 选或不选，容量不太大 & T301 & DFS 子集 \\
B13 基础数学 & 枚举因子到 n & 质因数、幂、约数 & T511/T512 & 小范围试除 \\
B14 矩阵优化 & 按题意多重循环超时 & 调整乘法/扫描顺序 & 第 6 章矩阵 & 暴力部分分 \\
B15 时间区间 & 每分钟逐个标记 & 时间范围不大或可差分 & T103/T110 & 数组标记小范围 \\
\bottomrule
\end{longtable}
\endgroup

\section{C 题快速判断树}

C 题最怕一边读一边写。正确流程是先拆对象、状态、命令、规则。

\subsection{C 题读题分解}

\begin{longtable}{p{3.0cm}p{6.0cm}p{5.0cm}}
\toprule
要找什么 & 在题面中怎么标记 & 代码落点 \\
\midrule
对象 & 用户、文件、节点、任务、请求、设备、词条 & struct \\
状态 & 空闲/占用、有效/过期、允许/拒绝、在线/离线 & enum 或 int \\
命令 & create、delete、query、request、release & T602 命令分发 \\
时间 & 到期、当前时刻、持续时间、周期 & T110、T514 \\
规则优先级 & 先判断 A，再判断 B；相同时间谁优先 & 独立函数，写注释 \\
输出 & 每次命令输出、最后统一输出、格式化文本 & 先写输出样例函数 \\
\bottomrule
\end{longtable}

\subsection{C 题模板选择}

\begin{longtable}{p{4.6cm}p{5.4cm}p{4.0cm}}
\toprule
题面信号 & 做法 & 模板 \\
\midrule
长题面、多对象、多规则 & 大模拟拆解 & T601 \\
每行一个命令 & 命令分发 & T602 \\
对象状态变化 & 状态机 & T603 \\
按时间发生 & 事件排序 & T110 \\
每次取最优任务 & 优先队列 / 懒删除 & T604 \\
路径、URL、配置 & 字符串解析 & T501 / T606 / T607 \\
括号、表达式、方程式 & 栈 / 递归解析 / 扫描 & T608 / T609 \\
缓存淘汰 & list + map 或队列 & T605 \\
地图带状态 & 状态 BFS & C 题章节 \\
树或目录结构 & DFS、父子 map、路径查找 & T606 / T208 \\
\bottomrule
\end{longtable}

\subsection{C 题详细叶子表}

\begingroup
\small
\setlength{\tabcolsep}{3pt}
\begin{longtable}{p{2.7cm}p{3.2cm}p{3.2cm}p{3.0cm}p{2.6cm}}
\toprule
叶子 & 题面信号 & 代码结构 & 易错点 & 模板 \\
\midrule
C1 命令分发 & 每行操作名不同 & 主循环读 op，分函数处理 & 参数数量不同 & T602 \\
C2 状态机 & 对象有状态变化 & 状态字段 + 转移函数 & 同一操作多次转移 & T603 \\
C3 事件排序 & 时间、到期、归还、同时发生 & Event 结构体排序 & 同时事件优先级 & T110 \\
C4 优先队列 & 每次取最优任务/对象 & heap 或 set & 失效元素删除 & T604 \\
C5 路径目录 & 文件、目录、父子、配额 & 树节点 + path map & 创建失败回滚 & T606 \\
C6 URL/参数 & 路径、查询参数、匹配规则 & split + 参数 map & 空参数、类型检查 & T607 \\
C7 JSON/配置 & 嵌套键值、查询路径 & 栈/递归建结构 & 转义和层级 & T501 \\
C8 Markdown/文本 & 多行文本、块、行内格式 & 先分块再处理行 & 空行分块 & T601 \\
C9 表达式/括号 & 运算、括号、化学式 & 栈或递归解析 & 优先级、嵌套倍数 & T608/T609 \\
C10 编码解码 & 十六进制、压缩、校验 & 指针扫描字符串 & 长度字段越界 & T513 \\
C11 矩阵图像 & 固定矩阵、扫描顺序、变换 & 二维数组 + 顺序表 & 行列、取整 & T102 \\
C12 权限规则 & 用户、角色、资源、权限 & set/map 规则判断 & 通配符、最大等级 & T603 \\
C13 带状态 BFS & 地图、钥匙、方向、时间 & visited 包含状态 & 状态维度漏掉 & T201 \\
C14 树/目录模拟 & 子树、删除、查询、合并 & DFS + 父子关系 & 删除后引用失效 & T208 \\
C15 分层拿分 & 规则太多写不完 & 先实现核心命令 & 输出格式别乱 & T601 \\
\bottomrule
\end{longtable}
\endgroup

\section{逐句读题分支树}

考场读题时，建议拿笔在题面旁边标记下面六类词。标完后再决定模板。

\begin{longtable}{p{2.8cm}p{5.8cm}p{5.4cm}}
\toprule
标记 & 看到什么就圈出来 & 后续判断 \\
\midrule
对象 & 学生、用户、文件、节点、仓库、服务器、任务、词条 & 是否需要 struct？是否有编号到对象的 map？ \\
属性 & 权重、时间、状态、坐标、容量、权限、父节点 & 是否要排序、比较器、状态机？ \\
动作 & 查询、修改、创建、删除、移动、合并、授权、到期 & 是否是命令分发或数据结构？ \\
范围 & $n$、$m$、$q$、值域、字符串总长、矩阵大小 & 直接排除不可能复杂度。 \\
关系 & 相邻、包含、依赖、祖先、连通、路径、前缀 & 图、树、拓扑、字符串、区间？ \\
目标 & 最大、最小、数量、是否存在、方案数、输出过程 & 统计、二分、DP、搜索还是模拟？ \\
\bottomrule
\end{longtable}

\subsection{从句子到模板}

\begin{longtable}{p{4.8cm}p{5.2cm}p{4.2cm}}
\toprule
题面句子模式 & 立即追问 & 模板方向 \\
\midrule
“给定 $q$ 次询问” & 每次询问是否只读不改？ & 前缀和 / 数据结构 \\
“每次修改后回答” & 修改是点、区间还是对象状态？ & 树状数组 / 线段树 / 状态机 \\
“求满足条件的数量” & 条件能否拆成前缀、排序或哈希？ & 计数 / 排序扫描 / map \\
“求最小的 $x$ 使得……” & $x$ 越大是否越容易满足？ & 二分答案 \\
“选择若干个” & 每个物品选一次还是多次？ & 01 背包 / 完全背包 \\
“从 A 到 B 的最短……” & 边权是否全为 1？ & BFS / Dijkstra \\
“任务必须在前置任务后” & 是否有环？ & 拓扑排序 \\
“字符串由若干规则组成” & 是否有括号或嵌套？ & split / 栈 / 递归解析 \\
“路径包含 . 和 ..” & 是否需要规范化？ & T606 \\
“状态在某时刻过期” & 过期是在操作前还是操作后处理？ & 状态机 / 事件模拟 \\
\bottomrule
\end{longtable}

\section{B 题专项分支树}

B 题常见陷阱是“暴力能写但会超时”。先写出暴力，再问它重复算了什么。

\subsection{B 题从暴力到优化}

\begin{longtable}{p{4.4cm}p{5.2cm}p{4.4cm}}
\toprule
暴力形态 & 重复在哪里 & 优化方向 \\
\midrule
每次询问都扫一段数组 & 区间和/数量重复计算 & 一维前缀和 T101 \\
每次询问都扫一个矩形 & 子矩阵重复计算 & 二维前缀和 T102 \\
每次区间加都逐个改 & 区间修改重复 & 差分 T103/T104 \\
每次判断点是否存在都扫所有点 & 存在性查询重复 & set / unordered\_set \\
每次找匹配点都双重循环 & 坐标或编号匹配重复 & map / 双指针 / 排序 \\
每个候选答案都完整模拟 & 答案范围大 & 二分答案 T107 \\
每次重新排序 & 只需维护极值或有序集合 & priority\_queue / set \\
每次从头判断依赖任务 & 依赖结构固定 & 拓扑排序 T205 \\
\bottomrule
\end{longtable}

\subsection{B 题二分答案确认表}

\begin{longtable}{p{4.6cm}p{4.8cm}p{4.6cm}}
\toprule
问题 & 是 & 否 \\
\midrule
答案是一个整数或实数吗？ & 继续问单调性 & 多半不是二分答案。 \\
给定答案 $x$，能快速判断可不可行吗？ & 写 \Code{check(x)} & 不适合二分。 \\
$x$ 变大后，可行性只朝一个方向变化吗？ & 可以二分 & 不能二分。 \\
题目问最小时间/最大阈值/最少资源吗？ & 优先尝试二分 & 可能是 DP/贪心。 \\
\bottomrule
\end{longtable}

\section{C 题专项分支树}

C 题不是先想算法，而是先建模。建模错误比模板不会更致命。

\subsection{C 题对象建模表}

\begin{longtable}{p{3.6cm}p{5.2cm}p{5.0cm}}
\toprule
题面对象 & 常见字段 & 容器选择 \\
\midrule
用户 / 角色 / 权限 & 名称、权限集合、等级、资源类型 & \Code{map<string, Obj>}，\Code{set<string>} \\
文件 / 目录 & 名称、父节点、子节点、大小、配额 & map 路径到节点，或树节点数组 \\
服务器 / 节点 & 编号、算力、任务数、可用状态 & vector + set/priority\_queue \\
任务 / 请求 & 时间、优先级、资源需求、状态 & struct + sort/heap \\
词条 / token & 字符串、权重、出现位置 & map/unordered\_map + heap \\
地图状态 & 坐标、方向、钥匙、时间、能量 & queue + visited 多维数组/map \\
\bottomrule
\end{longtable}

\subsection{C 题规则优先级表}

\begin{longtable}{p{4.6cm}p{5.0cm}p{4.6cm}}
\toprule
题面出现 & 必须单独确认 & 常见错误 \\
\midrule
“若同时发生” & 同时事件的排序规则 & 先后顺序写反。 \\
“如果不存在则……” & 缺省值或失败输出 & map 自动创建误判存在。 \\
“到期” & 到期时刻是否还有效 & $<$ 和 $\le$ 写错。 \\
“优先选择” & 比较器所有关键字 & 忘记编号升序等最后规则。 \\
“忽略非法操作” & 是否改变状态、是否输出 & 非法时仍修改了数据。 \\
“回滚/失败” & 前面临时改动要撤销 & 配额、路径创建类常见。 \\
\bottomrule
\end{longtable}

\subsection{C 题字符串解析分支}

\begin{longtable}{p{4.6cm}p{5.2cm}p{4.2cm}}
\toprule
字符串形态 & 判断 & 模板 \\
\midrule
固定分隔符，无嵌套 & split 即可 & T501 \\
路径，含 \Code{.}、\Code{..}、多个斜杠 & 栈式规范化 & T606 \\
URL，含路径和查询参数 & 先按 \Code{?}，再按 \Code{\&} & T607 \\
括号嵌套 & 栈或递归 & T608 \\
有运算优先级 & 两栈或递归下降 & T609 \\
化学式/表达式含倍数 & 递归解析返回 map & C 题表达式扩展 \\
多行文本块 & 先分块，再行内解析 & T601 \\
\bottomrule
\end{longtable}

\section{D/E 题快速判断树}

D/E 题不要先问“我会不会满分”，先问“我能不能拿部分分”。

\begin{enumerate}
  \item 找子任务数据范围：$n \le 20$、链、树、无修改、边权 1、查询少。
  \item 写最朴素暴力，估算它能过哪档。
  \item 如果暴力只有 0 分，再找一个最接近的模板：BFS、Dijkstra、Kruskal、树状数组、线段树、DP。
  \item 优先实现你能调通的版本，不把时间全部压在复杂满分做法上。
\end{enumerate}

\begin{longtable}{p{4.6cm}p{5.2cm}p{4.2cm}}
\toprule
题面信号 & 满分方向 & 部分分方向 \\
\midrule
$n \le 20$ 子任务 & 状压 / DFS & 枚举子集 \\
树上路径 & LCA / 树上差分 & 每次 DFS 找路径 \\
区间修改查询 & 线段树 & 暴力 / 树状数组 / 差分 \\
最小生成代价 & Kruskal & 小图枚举或直接 Kruskal \\
字符串匹配 & KMP / Trie & find / 暴力匹配 \\
正权最短路 & Dijkstra & 小图 Floyd / 无权 BFS \\
最大化最小值 & 二分答案 & 枚举答案小范围 \\
复杂 DP & 优化 DP & DFS / 小状态 DP \\
动态集合取最值 & 平衡树 / 堆懒删除 & 每次排序或线性扫描 \\
\bottomrule
\end{longtable}

\section{D/E 部分分决策树}

下面这棵树只服务一件事：不会满分时，尽快产出能提交的代码。

\begin{enumerate}
  \item 先看子任务表，有没有 $n \le 20$、$n \le 100$、链、树、无修改、边权 1、查询少。
  \item 如果有小 $n$：优先 DFS、枚举、Floyd、$O(n^2)$ DP。
  \item 如果有特殊图：边权 1 用 BFS，树用 DFS，DAG 用拓扑。
  \item 如果有特殊操作：无修改用前缀，有修改少就每次重算，查询少就每次暴力。
  \item 如果只有大数据：选择你最熟的相近模板，写 30 到 50 分版本。
\end{enumerate}

\begin{longtable}{p{4.4cm}p{5.2cm}p{4.4cm}}
\toprule
子任务/特殊性质 & 可写版本 & 常用模板 \\
\midrule
$n \le 20$ & 枚举子集、DFS、状压 & 状压 DP、搜索 \\
$n \le 100$ & Floyd、$O(n^3)$ DP、暴力查询 & T206 \\
树退化成链 & 数组/前缀处理路径 & T101 / T103 \\
没有修改 & 预处理后回答查询 & 前缀、DFS 序 \\
查询次数很少 & 每次 DFS/BFS/扫描 & 暴力部分分 \\
边权全为 1 & BFS 替代 Dijkstra & T201 \\
图是 DAG & 拓扑 DP & T205 \\
值域小 & 计数数组/桶 & T106 \\
参数 $k$ 很小 & 状态压缩或按 $k$ 枚举 & 状压入口 \\
\bottomrule
\end{longtable}

\subsection{D/E 题详细叶子表}

\begingroup
\small
\setlength{\tabcolsep}{3pt}
\begin{longtable}{p{2.6cm}p{3.2cm}p{3.2cm}p{3.0cm}p{2.8cm}}
\toprule
叶子 & 识别信号 & 满分方向 & 部分分写法 & 翻到哪里 \\
\midrule
DE1 小 n & $n \le 20$ 或参数小 & 状压/搜索优化 & 枚举子集、DFS & 第 8 章状压 \\
DE2 小图 & $n \le 100$ & Floyd/DP & Floyd 或暴力建图 & T206 \\
DE3 树路径 & 树上多次路径 & LCA/树剖/树差分 & 每次 DFS 找路径 & T208 \\
DE4 子树查询 & 子树加、子树查 & DFS 序 + 树状数组/线段树 & 每次遍历子树 & T401/T402 \\
DE5 动态区间 & 区间改查交替 & 线段树 & 小范围数组模拟 & T402 \\
DE6 动态集合 & 插入删除查极值 & set/堆懒删除 & 每次排序/扫描 & T604 \\
DE7 最短路 & 正权图路径 & Dijkstra/建图优化 & Floyd 小图、BFS 边权1 & T204 \\
DE8 最小生成 & 连通最小代价 & Kruskal/Prim & 直接 Kruskal & T207 \\
DE9 字符串匹配 & 长串、多模式、前缀 & KMP/Trie/AC 自动机 & find 暴力 & T502/T503 \\
DE10 复杂 DP & 最优方案数、限制多 & 状态设计/优化 & 记忆化 DFS 小数据 & 第 8 章 DP \\
DE11 数学计数 & 组合、取模、约数 & 预处理/数论 & 枚举小范围 & T511/T512 \\
DE12 离线处理 & 操作可重排、查询很多 & 排序/分治/离线 DSU & 按输入暴力 & T110/T203 \\
DE13 几何 & 点线面、距离、交叉 & 计算几何 & 小数据枚举 & 后续补强 \\
DE14 网络流/匹配 & 分配、匹配、容量流 & 网络流/二分图 & 贪心或小图 DFS & 放弃满分 \\
\bottomrule
\end{longtable}
\endgroup

\section{容易误判的分支}

\begin{longtable}{p{4.5cm}p{4.8cm}p{4.8cm}}
\toprule
看起来像 & 但如果题面出现 & 应改判为 \\
\midrule
简单模拟 & $q$ 很大，每次都查区间 & 前缀和 / 差分 / 数据结构 \\
排序题 & 要动态插入删除后仍排名 & set、树状数组、线段树 \\
字符串题 & 有括号嵌套、表达式优先级 & 栈或递归下降 \\
图搜索题 & 边权不是 1 & Dijkstra，不是普通 BFS \\
BFS 题 & 状态包含钥匙、方向、时间、能量 & 状态 BFS \\
贪心题 & 选或不选、容量限制明显 & 背包 DP \\
二分题 & check 不单调 & 不能二分，改 DP/贪心/扫描 \\
树题 & 只有父子目录和路径字符串 & 可能是大模拟，不一定是 LCA \\
数学题 & 有多次询问和大范围 & 预处理、筛法、快速幂 \\
\bottomrule
\end{longtable}

\section{无法判断时的保底流程}

如果读完题仍不知道算法：

\begin{enumerate}
  \item 写出最暴力做法。
  \item 估算暴力复杂度。
  \item 看数据范围中有没有小数据能过。
  \item 找特殊性质：无修改、链、树、边权 1、参数小。
  \item 先实现能过样例和小数据的版本。
  \item 再尝试从本章关键词表找优化方向。
\end{enumerate}

\section{判断树反向测试表}

刷完一道题后，用这张表检查第 3 章是否真的帮你判断。如果某行填不上，就说明手册还要继续补。

\begin{longtable}{p{4.0cm}p{5.0cm}p{5.0cm}}
\toprule
问题 & 例子 & 应能定位到 \\
\midrule
题面主体是什么？ & 数组、矩阵、图、字符串、对象、事件 & 题面结构第一判断 \\
是否有多次操作？ & $q$ 次查询、每次修改、每次输出 & 操作模式判断树 \\
暴力复杂度是多少？ & $O(nq)$、$O(n^2)$、$O(nm)$ & 数据范围判断树 \\
关键词是否误导？ & 看起来像模拟但 $q$ 很大 & 容易误判分支 \\
满分不会时能拿什么？ & 小数据、特殊性质、暴力 & D/E 部分分决策树 \\
对应模板编号是什么？ & T101、T107、T603 等 & 关键词到模板总索引 \\
\bottomrule
\end{longtable}

\section{最后 5 分钟判断}

\begin{longtable}{p{4.2cm}p{9.8cm}}
\toprule
情况 & 处理 \\
\midrule
A/B 还没过样例 & 放下后面题，先修 A/B。 \\
C 写了一半但样例不过 & 注释掉高风险优化，保留最基础模拟。 \\
D/E 没思路 & 找最小子任务，写暴力或特殊性质。 \\
数组越界风险高 & 全部下标检查一遍，必要时数组开大 5 到 10。 \\
输出格式不确定 & 对照样例逐字符检查空格、换行、小数位。 \\
\bottomrule
\end{longtable}
